\chapter*{Cómo leer esta guía}
\addcontentsline{toc}{chapter}{Cómo leer esta guía}
\markboth{Cómo leer esta guía}{}

Esta guía cuenta el proyecto \textbf{Torre del Caribe} de principio a fin para alguien que no estuvo cuando se hizo: una
integrante del equipo, un profesor que pregunta ``¿y esto de dónde salió?'' o el propio autor, dentro de un mes, cuando
ya no recuerde los detalles. Es la versión \textbf{sin fórmulas difíciles} de la \emph{Guía técnica}. Dice lo mismo y con
las mismas cifras, pero explica cada paso con palabras de todos los días, con comparaciones y con cuentas que se pueden
hacer en una calculadora.

\section*{Tres reglas de esta guía}
\begin{enumerate}
\item \textbf{Cada cifra es real.} Todos los números salen de fuentes oficiales (INEGI, SECTUR, el gobierno de Quintana
  Roo, la NOAA) o de un cálculo del proyecto que se puede repetir. Una prueba automática (\codigo{tests/test_guias.py})
  revisa que las cifras clave de esta guía sigan coincidiendo con los datos.
\item \textbf{Cada paso dice por qué.} No basta con decir ``se usó tal cosa''. Siempre se explica qué otras opciones había y
  por qué se eligió esa.
\item \textbf{Lo que no se sabe, se dice.} Cuando un dato no existe, la guía lo llama \emph{hueco} y no lo rellena.
\end{enumerate}

\section*{Las cajas de colores}
\begin{intuicion}
La idea de cada paso, explicada con una comparación de la vida diaria.
\end{intuicion}
\begin{ejemplo}
Una cuenta hecha con números reales del proyecto, paso a paso, para ver que no hay magia.
\end{ejemplo}
\begin{porque}
Las otras opciones que había, y la razón por la que se eligió esta.
\end{porque}
\begin{teprofe}
Una pregunta que probablemente hará un profesor en el coloquio, con una respuesta corta y defendible.
\end{teprofe}
\begin{error}
Algo que salió mal durante el proyecto, cómo se descubrió y cómo se corrigió. Estas cajas son importantes: muestran que
las cifras se revisaron y no se aceptaron a ciegas.
\end{error}

\section*{Cómo está organizada}
\begin{itemize}
\item \textbf{Parte I, La idea} (capítulos 1 a 3): el problema, las preguntas incómodas (¿por qué aparece Claude?, ¿por qué
  tantos archivos?) y de dónde salen los datos.
\item \textbf{Parte II, Cómo se hizo} (capítulos 4 a 11): una materia por capítulo, en el orden en que los datos viajan por
  el proyecto: guardarlos, encontrar patrones, predecir, medir la incertidumbre, repartir el dinero, vigilar cada semana,
  diseñar la campaña y mostrarla en la página.
\item \textbf{Parte III, Para defenderlo} (capítulos 12 y 13): lo que no se pudo medir, los errores que se corrigieron y
  cómo explicar el proyecto en uno y en cinco minutos.
\end{itemize}

Si se quiere la fórmula exacta de algo, el mismo tema está en la \emph{Guía técnica} con el mismo nombre. Si se quiere
saber quién decidió qué y cuándo, está en \emph{Las decisiones, fase por fase}.
